\documentclass[journal]{IEEEtran}

\usepackage[utf8]{inputenc}
\usepackage[T1]{fontenc}
\usepackage{amsmath,amssymb,amsfonts}
\usepackage{graphicx}
\usepackage{booktabs}
\usepackage{array}
\usepackage{url}
\usepackage[hidelinks]{hyperref}
\usepackage{orcidlink}

\begin{document}

\title{Horizon-Dependent Value of Geodetic Context in Operational Earthquake
Forecasting: A Leakage-Free, Multi-Region Study with Pre-Registered Negative Results}

\author{%
  Felipe Santib\'a\~nez-Leal~\orcidlink{0000-0002-0150-3246},~\IEEEmembership{Independent Researcher}%
  \thanks{F.\ Santib\'a\~nez-Leal is with the CAOS open-research programme,
    Santiago, Chile (e-mail: \texttt{fsantibanezleal@ug.uchile.cl};
    ORCID: 0000-0002-0150-3246).}%
  \thanks{Preprint, 2026. Code, configurations, provenance manifests, and the
    committed daily/back-analysis artifacts are released under MIT at
    \url{https://github.com/fsantibanezleal/CAOS_SEISMIC}; a static forecast
    viewer is live at \url{https://seismic.fasl-work.com}. This is an
    independent research and education tool; it is not an alarm system and must
    not be used for life-safety decisions.}%
}

\markboth{Preprint, 2026 --- CAOS open-research programme}{Santib\'a\~nez-Leal:
Horizon-Dependent Value of Geodetic Context in Operational Earthquake Forecasting}

\IEEEtitleabstractindextext{%
\begin{abstract}
Operational earthquake forecasting (OEF) issues calibrated conditional
probabilities of future seismicity; the Epidemic-Type Aftershock Sequence (ETAS)
model is its de-facto benchmark, and no machine-learning temporal point process
has been shown to beat a well-fit ETAS prospectively under fair testing. We build
a leakage-free, multi-region OEF testbed on a global moment-magnitude catalog
($M \ge 5$, 1990--2026), score every model with the community-standard CSEP
framework (pyCSEP) in information gain per earthquake (IGPE, nats), and pre-register
ship/no-ship rules with shuffled-label negative controls. Within it we report four
findings, each artifact-backed and honestly framed. \emph{First}, the value of a
GNSS-strain geodetic context covariate, added to a Hawkes-structured neural point
process, is \emph{horizon-dependent}: at the 1--7 day operational horizon the
strain-conditioned model does not beat tiled ETAS (mean IGPE $-0.053$ over eight
weekly windows, only $4/8$ positive), whereas at a 30 day horizon it beats ETAS
robustly (global mean IGPE $+0.115$ over 2166 events, $9/10$ windows positive,
paired-$t$ lower bound $+0.087$; positive in all four high-seismicity regions).
The mechanism is explicit: the calibrated model deploys a time-flat geodetic
\emph{background}, not a triggering model, so it wins where the background
dominates the integral (30 days) and loses where ETAS triggering dominates
(1--7 days). \emph{Second}, a convex log-score-optimal stack of ETAS-family
variants earns a small but significant global gain ($+0.011$ over 751 events) that
does \emph{not} generalize to the canonical active margins; under the
pre-registered rule it is a no-ship, and a temporally adaptive variant is refuted
as a multiple-comparison artifact. \emph{Third}, ETAS adds skill over a stationary
smoothed-seismicity Poisson null only where there is triggering to exploit,
quantifying the high-versus-low-seismicity bias rather than asserting it.
\emph{Fourth}, the binding 1--7 day consistency failure is count over-dispersion,
not spatial shape: a catalog-based (branching-simulation) number test passes a
window the Poisson number test rejects, and the frozen-intensity forecast
systematically under-counts by $\approx 28\%$ because it omits within-window
secondary triggering. We conclude that base tiled ETAS is at or near the practical
ceiling for mean-rate 1--7 day IGPE over global $M \ge 5$, and that the honest
place for a geodetic covariate is a longer, background-dominated outlook.
\end{abstract}

\begin{IEEEkeywords}
operational earthquake forecasting, ETAS, epidemic-type aftershock sequence,
neural temporal point process, geodetic strain, GNSS, CSEP, pyCSEP, information
gain, prospective evaluation, over-dispersion, reproducibility.
\end{IEEEkeywords}}

\maketitle
\IEEEdisplaynotcompsoctitleabstractindextext
\IEEEpeerreviewmaketitle

% =====================================================================
\section{Introduction}
% =====================================================================

\IEEEPARstart{E}{arthquakes} cannot be predicted deterministically: whether a
small rupture cascades into a great earthquake depends on unmeasurably fine
details of the crust~\cite{geller1997}. What \emph{can} be done, and is done
operationally, is \emph{forecasting}: issuing a probability strictly in $(0,1)$
for a region, a magnitude band, and a short horizon. Following the ICEF
definition~\cite{jordan2011}, a prediction is a deterministic statement that an
event will or will not occur, whereas a forecast reports a probability with
uncertainty. Short-term probabilities of a large event may vary over orders of
magnitude relative to background yet typically remain small in absolute terms
(well under a few percent per day)~\cite{jordan2011}. This distinction is not
pedantic: the documented harm in the field is \emph{false reassurance}, and every
number here is a conditional probability shown next to its long-term baseline,
scored prospectively, and never rendered as an alarm.

The de-facto short-term baseline is the Epidemic-Type Aftershock Sequence (ETAS)
model~\cite{ogata1988,ogata1998}, a self-exciting Hawkes point process whose
conditional intensity superposes a stationary background on the summed,
Omori--Utsu-decaying, Gutenberg--Richter-weighted offspring of every past event.
Operational systems such as the USGS aftershock forecast~\cite{reasenberg1989,page2016}
and OEF-Italy~\cite{spassiani2023} run ETAS-class models as scheduled services and
publish calibrated probabilities with uncertainty. A decade of prospective
next-day forecasting over California confirms that no single model dominates: the
short-term-clustering models excel during aftershock sequences, and ETAS is the
consistent generalist~\cite{serafini2025}.

Against this baseline, machine learning has repeatedly promised more and, under
fair testing, repeatedly failed to deliver a prospective gain. The EarthquakeNPP
benchmark~\cite{stockman2026} evaluated five modern neural spatio-temporal point
processes on California with strict chronological splits and CSEP consistency
tests, and found that \emph{none} outperformed ETAS, precisely on the spatial test
where forecasting value lives; the apparent advantages of earlier neural work
trace to a data-leakage flaw (non-chronological splits) and to excluding the large
sequences. The canonical cautionary tale is a deep network for aftershock
patterns~\cite{devries2018} whose reported skill was matched by a two-parameter
logistic regression on a single feature~\cite{mignan2019}, with the area under the
ROC curve exposed as blind to the calibration a forecast must publish. Where
learned models \emph{can} help ETAS, the gains come from two real gaps, not from
network depth: ingesting covariates ETAS cannot easily absorb, and learning
spatial anisotropy~\cite{dascher2023,zlydenko2023}.

This paper studies one such covariate, \emph{geodetic strain from GNSS}, and asks
not whether it helps but \emph{at which horizon}. Geodesy is an established input
to long-term global rate models: the GEAR1 model blends smoothed seismicity with a
strain-rate model~\cite{bird2015}, and prospective global hybrids find geodetic
strain adds skill globally while its regional contribution is
weaker~\cite{bayona2022}. We port that intuition into a short-term OEF testbed and
measure it prospectively. Our contributions are:
\begin{enumerate}
\item A leakage-free, multi-region OEF testbed with a strict forecast clock,
  CSEP/pyCSEP scoring in IGPE, pre-registered ship rules, and shuffled-label
  negative controls (Sec.~\ref{sec:methods}).
\item The \emph{horizon-dependent} value of a GNSS-strain context covariate: no
  gain at 1--7 days, a robust and regionally consistent gain at 30 days, with a
  mechanistic explanation (Sec.~\ref{sec:geodetic}).
\item Pre-registered \emph{negative} results: ETAS-family stacking that does not
  generalize across regions, and the 1--7 day mean-rate IGPE ceiling
  (Sec.~\ref{sec:stack}).
\item A methodological result: over-dispersion-honest, catalog-based consistency
  scoring, and the $\approx 28\%$ secondary-cascade under-count of a
  frozen-intensity forecast (Sec.~\ref{sec:overdisp}).
\end{enumerate}
Every result is backed by a committed, reproducible artifact; negative and null
results are reported, not hidden.

% =====================================================================
\section{Data}
\label{sec:data}
% =====================================================================

The catalog spine is USGS ComCat (FDSN event service, public domain), homogenized
to moment magnitude $M_w$. A per-network station-pair regression to $M_w$ is not
always available, so the Scordilis~\cite{scordilis2006} global $m_b \to M_w$ and
$M_s \to M_w$ orthogonal-regression relations are used as a documented fallback,
with per-event provenance recorded; this recovers a $6\times$ larger conditioning
catalog (from $1.65\times10^4$ to $9.7\times10^4$ events, $M \ge 4.5$, 1990--2026)
at the cost of a bounded, audited $b$-value inflation that the per-tile fit largely
removes (naive global $b \approx 1.59$; per-tile production $b \approx 1.34$).
The magnitude of completeness $M_c$ and the Gutenberg--Richter $b$-value are
estimated on rolling windows~\cite{wiemer2000,woessner2005}. Declustering, where
needed for the stationary background, uses stochastic declustering~\cite{zhuang2002}
and the nearest-neighbour method~\cite{zaliapin2013}. The geodetic covariate is the
GNSS strain rate from the Nevada Geodetic Laboratory MIDAS velocity
field~\cite{blewitt2016} (20{,}168 stations after an endpoint and column-parse fix),
gridded to a coarse global field; strain is high on active margins (Japan 95,
California 154, Chile 64 nstrain/yr) and low in stable interiors (Central US 3.6,
Western Europe 5.6 nstrain/yr). Plate-boundary geometry is PB2002~\cite{bird2003}.
Raw catalogs, features, and model weights are rebuildable from the code and
provenance manifests and are not versioned; the compact daily and back-analysis
artifacts are.

% =====================================================================
\section{Methods}
\label{sec:methods}
% =====================================================================

\subsection{Baseline models}

\textbf{ETAS.} The space--time conditional intensity~\cite{ogata1998} is
\begin{equation}
\lambda(t,x,y \mid \mathcal{H}_t) = \mu(x,y) + \!\!\sum_{i:\,t_i < t}\!\! K\,
e^{\alpha(M_i - M_0)} \Big(1 + \tfrac{t-t_i}{c}\Big)^{-p}\, f(x-x_i, y-y_i \mid M_i),
\end{equation}
with the Ogata inverse-power spatial kernel
$f(x,y \mid M_i) = \frac{q-1}{\pi \zeta^2}\big(1 + r^2/\zeta^2\big)^{-q}$,
$\zeta = D\,e^{\gamma(M_i-M_0)}$. Parameters are fit by maximizing the
point-process log-likelihood
\begin{equation}
\ln L = \sum_i \ln \lambda(t_i,x_i,y_i \mid \mathcal{H}_{t_i})
- \int_0^T\!\!\int_A \lambda(t,x,y\mid\mathcal{H}_t)\,dx\,dy\,dt,
\label{eq:ppll}
\end{equation}
subject to two stability gates: finite branching ($\alpha < \beta$, $\beta = b\ln 10$)
and subcriticality (branching ratio $n < 1$). Because a single global ETAS over a
$10^5$-event catalog is both intractable and physically wrong (subduction and
stable interior obey different kinematics), we fit a \emph{regime-tiled} ETAS:
195 ETAS tiles are fit and 265 sparse tiles fall back to their smoothed null, on
$M_c = 5.35$, $b = 1.337$.

\textbf{Smoothed-seismicity Poisson null.} The stationary background
$\mu(x,y) = \sum_i K_{d_i}(r)$ uses the Helmstetter--Kagan--Jackson adaptive
power-law kernel with bandwidth set by the distance to the $n$-th nearest
neighbour~\cite{helmstetter2007}. It serves both as the ETAS background field and,
crucially, as the \emph{mandatory stationary null} any time-dependent model must
beat. \textbf{Reasenberg--Jones}~\cite{reasenberg1989,page2016} is a transparent
aftershock-only fallback and sanity check.

\subsection{Neural challenger with geodetic context}

The challenger is a context-conditioned temporal point process
(\texttt{context\_tpp}): a CNN encodes a gridded covariate field (here the
\texttt{gnss\_strain\_rate} channel; other geophysical channels are honestly
zero-filled) into a Hawkes-structured conditional intensity, in the spirit of
RECAST~\cite{dascher2023} and FERN~\cite{zlydenko2023} (keep the additive
background plus summed-triggering skeleton, replace fixed kernels with learned
maps). It is trained on the global catalog on a single laptop GPU. A fit-time
scalar \texttt{rate\_cal} anchors the integrated one-day forecast rate to the
training rate on the same multi-resolution grid the evaluation uses; without it the
background head over-predicts the absolute rate by $\approx 490\times$ (a failure we
diagnose and report, Sec.~\ref{sec:geodetic}). The challenger is a \emph{gated}
model: it reaches the public map only if it beats ETAS prospectively \emph{and} is
calibrated; until then, ETAS ships.

\subsection{Evaluation: CSEP, IGPE, and the forecast clock}

Skill is established with the community CSEP framework via
pyCSEP~\cite{savran2022}. Consistency (calibration) is tested with the
number (N), magnitude (M), spatial (S), and conditional-likelihood (CL)
tests~\cite{schorlemmer2007,zechar2010}, built on the Poisson joint log-likelihood
over space--magnitude bins. \emph{Passing consistency is necessary but not
sufficient}: skill requires \emph{winning a comparison test} against both the
Poisson null and ETAS, measured as information gain per earthquake (IGPE, nats):
\begin{equation}
I_N(A,B) = \frac{1}{N}\sum_{i=1}^N\!\big(\ln\lambda_A(k_i) - \ln\lambda_B(k_i)\big)
- \frac{\hat{N}_A - \hat{N}_B}{N},
\label{eq:igpe}
\end{equation}
tested with the paired $t$-test $T = I_N/(s/\sqrt{N}) \sim t_{N-1}$ and corroborated
by the Wilcoxon signed-rank W-test~\cite{rhoades2011}. ROC/AUC is deliberately
\emph{banned} as a skill metric: it is invariant to calibration and, on rare
per-cell-per-day tasks, degenerates into a region classifier (the DeVries
trap)~\cite{devries2018,mignan2019}. Proper scoring rules~\cite{gneiting2007} and a
Molchan/area-skill view~\cite{zechar2008} are reported only as communication aids.

\textbf{Leakage is structural.} A strict forecast clock hands the model only the
catalog slice up to issue time; each forecast is logged immutably with its exact
input state, so any past forecast is byte-reproducible and scored against the
catalog as it was at issue time. The authoritative measure is a
\emph{pseudo-prospective back-analysis}: fit the model family at the earliest
cutoff, then \emph{recondition} forward across weekly, non-overlapping,
fully-scorable windows (refresh the parents, hold the parameters and background),
scoring IGPE per window and per regional view. Reconditioning reproduces
full-refit-every-day to $0.05\%$ median deviation while running $13\times$ faster.
Ship/no-ship rules are \emph{pre-registered} before each run, and a shuffled-label
negative control guards against over-fitting.

% =====================================================================
\section{Results}
% =====================================================================

\subsection{ETAS adds skill only where there is triggering to exploit}
\label{sec:bias}

The first adversarial question is whether ETAS only looks good because it over-fits
high-seismicity margins. A 60-day, leakage-free, per-country back-analysis
(1-day horizon, IGPE vs.\ the Poisson null) answers it quantitatively
(Table~\ref{tab:bias}): ETAS adds skill in the densest triggering regime (Japan,
$+0.072$) and collapses to the null in low-seismicity interiors (Central US,
Western Europe, Eastern Australia, all $0.0$). A single-window benchmark tells the
same story at larger absolute scale (global $+3.71$ nats vs.\ the null), and shows
that Reasenberg--Jones alone is no better than the null. Low-seismicity $0.0$ is not
a failure; it is the correct statement that a self-exciting model has nothing to add
without aftershock sequences.

\begin{table}[!t]
\caption{ETAS IGPE vs.\ the Poisson null, 60-day leakage-free back-analysis.
Skill concentrates in the high-seismicity, triggering-rich regime.}
\label{tab:bias}
\centering
\begin{tabular}{llrr}
\toprule
View & Class & IGPE (1\,d) & IGPE (7\,d) \\
\midrule
Japan & high & $+0.072$ & $+0.062$ \\
Chile & high & $+0.0015$ & $+0.0024$ \\
California & high & $\sim 0$ & $+0.0016$ \\
New Zealand & high & $\sim 0$ & $-0.001$ \\
Central US & low & $0.0$ & $0.0$ \\
Western Europe & low & $0.0$ & $0.0$ \\
Eastern Australia & low & $0.0$ & $0.0$ \\
\bottomrule
\end{tabular}
\end{table}

\subsection{ETAS-family stacking: a global gain that does not generalize}
\label{sec:stack}

The only prospective lever the literature supports over a single well-fit ETAS is a
score-weighted ETAS-family ensemble, and only by a small, often non-significant
margin~\cite{marzocchi2012,herrmann2023}. A naive equal-weight pool of
heterogeneous models under-performs ETAS because averaging in the null dilutes the
triggering signal; so we build a convex, log-score-optimal stack of three
\emph{ETAS-family} variants (base tiled, short-memory, long-memory), with the
smoothed null entering only as each member's background, weights shrunk toward the
base vertex, and a pre-registered rule requiring a positive, CI-excludes-zero gain
on the global view \emph{and} in at least two high-seismicity regions at the 7-day
horizon.

The 14-window, 7-day back-analysis (Table~\ref{tab:stack}) earns a real, significant
\emph{global} gain ($+0.011$ over 751 events, $13/14$ windows positive, paired-$t$
CI excludes zero), carried by the long-memory variant capturing late-aftershock
structure the base cutoff misses. But only $1/4$ high-seismicity regions is positive
(California, on 3 events, noise). Under the pre-registered rule this is a
\textbf{no-ship}. A temporally adaptive variant appears to rescue it ($3/4$ regions
positive) but is refuted on scrutiny: it \emph{reduces} the global gain, and the
regional ``improvements'' are sign-flipping noise on 3--54 events (no regional CI
excludes zero), a textbook multiple-comparison artifact. The honest reading: the
regional 7-day $M \ge 5$ gate is under-powered; the global view is the only
adequately powered test, and there the static stack already gives the answer.

\begin{table}[!t]
\caption{Pre-registered multi-region gate for the ETAS-family stack, 7-day horizon,
14 weekly windows. Global gain is real; regional generalization fails the rule.}
\label{tab:stack}
\centering
\begin{tabular}{lrrl}
\toprule
View & mean IGPE vs.\ base & windows $+$ & paired-$t$ CI \\
\midrule
\textbf{Global} & $\mathbf{+0.0111}$ & $13/14$ & excludes zero $(+)$ \\
Chile & $-0.0001$ & $3/10$ & crosses zero \\
Japan & $-0.0008$ & $8/13$ & crosses zero \\
California & $+0.0077$ & $3/3$ (noise) & crosses zero \\
New Zealand & $-0.0008$ & $0/3$ & crosses zero \\
\bottomrule
\end{tabular}
\end{table}

\subsection{The 1--7 day failure is over-dispersion, not shape}
\label{sec:overdisp}

Why does the base ETAS N-test fail on active windows? Regional seismicity is
over-dispersed relative to Poisson (variance $\gg$ mean, from clustering), so the
Poisson number test over-rejects vigorous-but-plausible sequence counts. We add a
catalog-based number test~\cite{savran2020,kagan2017}: an ETAS branching-process
forward simulation produces the honest in-window $M \ge M_c$ count distribution.
On a 30-day global window (Table~\ref{tab:overdisp}), the frozen-intensity forecast
predicts $N_\text{fore} = 64.3$ (background 60.1 $+$ first-generation triggering
4.2), the full branching cascade lifts the mean to 82.1 with Fano factor 1.79, and
the observed 96 is consistent: the Poisson N-test fails at quantile $0.0001$, while
the catalog-based test passes at $0.133$. Two honest findings follow. The apparent
under-forecast is over-dispersion plus secondary cascade, not a rate bias. And the
frozen-intensity forecast \emph{systematically under-counts by $\approx 28\%$}
because it omits within-window secondary triggering; the catalog-based test is the
correct consistency tool. (This is one 30-day global window; the method generalizes.)

\begin{table}[!t]
\caption{Over-dispersion-honest consistency, 30-day global window. The
catalog-based test passes where the Poisson test over-rejects.}
\label{tab:overdisp}
\centering
\begin{tabular}{lr}
\toprule
Quantity & Value \\
\midrule
$N_\text{fore}$ (frozen intensity) & 64.3 \\
simulated mean (full cascade) & 82.1 $\;(+27.6\%)$ \\
Fano factor (Var/mean) & 1.79 \\
observed $M \ge M_c$ & 96 \\
Poisson N-test & $q = 0.0001$ \; (fail) \\
catalog-based N-test & $q = 0.133$ \; (pass) \\
\bottomrule
\end{tabular}
\end{table}

\subsection{The horizon-dependent value of geodetic context}
\label{sec:geodetic}

The headline result concerns the strain-conditioned neural challenger. Diagnosing
it required first fixing a calibration failure: the raw model over-predicted the
30-day global rate $\approx 490\times$ (98{,}932 vs.\ 248 observed) because the
training compensator constrained the background only at event locations, leaving the
background head free to inflate rate in event-free cells. The fit-time
\texttt{rate\_cal} anchor (computed on the evaluation grid) fixes this: the
calibrated model forecasts 218.9 events (ETAS 173.7, observed 248) and passes the
N-test. Only the calibrated model is scored below.

\textbf{At the 1--7 day operational horizon, the geodetic context does not beat
ETAS.} An 8-window, leakage-free, 7-day back-analysis gives mean IGPE vs.\ ETAS of
$-0.053$, positive in only $4/8$ windows, with high per-window variance (a single
badly-missed sequence window drives $-0.80$). The ETAS-family stack of
Sec.~\ref{sec:stack} likewise gains only $+0.009$ globally at 7 days and does not
generalize. The 1--7 day mean-rate IGPE is at or near a practical ceiling.

\textbf{At the 30-day horizon, the same geodetic context beats ETAS robustly and
regionally.} A 10-window back-analysis (Table~\ref{tab:geodetic}) gives a global
mean IGPE vs.\ ETAS of $+0.115$ over 2166 events, positive in $9/10$ windows, with a
paired-$t$ lower bound of $+0.087$ (excludes zero), and it is positive in
\emph{every} high-seismicity region. This is the opposite of the 7-day result.

\begin{table}[!t]
\caption{Strain-conditioned neural vs.\ ETAS, 30-day horizon, 10 leakage-free
windows. The geodetic-context gain is global \emph{and} regionally robust.}
\label{tab:geodetic}
\centering
\begin{tabular}{lrrr}
\toprule
View & mean IGPE vs.\ ETAS & windows $+$ & events \\
\midrule
\textbf{Global} & $\mathbf{+0.115}$ & $9/10$ & 2166 \\
Chile & $+0.656$ & $9/10$ & 47 \\
Japan & $+0.192$ & $8/10$ & 197 \\
California & $+1.007$ & $7/9$ & 11 \\
New Zealand & $+1.210$ & $8/8$ & 19 \\
\bottomrule
\end{tabular}
\end{table}

\textbf{The mechanism is explicit.} The calibrated model's total forecast is
essentially \emph{constant across windows} (18.1 events at 7 days, 77.3 at 30 days,
ratio $\approx 30/7$): it deploys a time-flat geodetic \emph{background}, not a
triggering model. That is precisely why it loses at 7 days (it cannot track evolving
sequences, where ETAS triggering dominates the integral) and wins at 30 days (a
better stationary background dominates the longer integral). This is consistent with
the global-scale geodetic literature, in which strain rate improves long-term global
rate models~\cite{bird2015} while its regional short-term contribution is
weaker~\cite{bayona2022}. The honest operational consequence is that the 1--7 day
product stays ETAS, and the geodetic context is deployed only as a longer,
background-dominated outlook where it measurably earns skill.

% =====================================================================
\section{Discussion}
% =====================================================================

The four results cohere into one message: for global $M \ge 5$ short-term
forecasting, base tiled ETAS is at or near the practical ceiling on mean-rate IGPE,
and the honest levers lie elsewhere. A separately conducted, adversarial review of
2024--2026 frontier paths (time-varying $b$, regime-gated neural routing, anisotropic
finite-fault kernels, magnitude-head networks, multi-resolution scoring) found that
none survived skeptical refutation as a large 1--7 day gain; the convergent verdict
matches the EarthquakeNPP conclusion~\cite{stockman2026} and our own measurements.
The value is not in chasing a larger 7-day number but in (i) scoring honestly
(catalog-based over-dispersion tests, Sec.~\ref{sec:overdisp}), (ii) deploying a
covariate at the horizon where it earns skill (Sec.~\ref{sec:geodetic}), and
(iii) refusing to ship a result that does not generalize (Sec.~\ref{sec:stack}).

The methodological posture is as much a contribution as any single number.
Pre-registered ship rules, shuffled-label negative controls, and the distinction
between a well-powered global test (751--2166 events) and an under-powered regional
gate (3--54 events) are what separate a real result from a multiple-comparison
artifact; the temporally adaptive stack, which a naive read would have shipped, is
the clearest example.

% =====================================================================
\section{Honest limitations}
% =====================================================================

The pseudo-prospective back-analysis spans weekly windows over a limited recent
period; the event counts are large (2166 global events at 30 days) but the temporal
span is short, so the results are best read as a well-powered snapshot rather than a
multi-year prospective record. The neural challenger's absolute calibration, though
fixed by \texttt{rate\_cal} for scoring, is not production-hardened, and it is not
used in the daily 1--7 day product. Only one context channel (GNSS strain) is
active; the remaining geophysical covariates are honestly zero-filled and not
claimed. The magnitude homogenization uses literature $M_w$ relations as a documented
fallback, with a bounded $b$-value bias that the per-tile fit largely removes. The
secondary-cascade under-count is measured on one 30-day global window. None of these
is hidden: each corresponds to a flag or an artifact in the released code.

% =====================================================================
\section{Conclusion}
% =====================================================================

We built a leakage-free, multi-region operational earthquake forecasting testbed
with pre-registered CSEP evaluation and used it to measure, rather than assert, four
things: that ETAS earns skill only where triggering exists; that an ETAS-family
stack gains globally but not regionally; that the binding 1--7 day consistency
failure is count over-dispersion with a $\approx 28\%$ secondary-cascade under-count;
and that the value of a GNSS-strain geodetic covariate is horizon-dependent, absent
at 1--7 days and robust at 30 days, because the calibrated model contributes a
time-flat background rather than triggering. The practical recommendation is to ship
calibrated ETAS at the operational horizon, score with catalog-based over-dispersion
tests, and reserve the geodetic context for a longer background-dominated outlook.
All code, configurations, provenance manifests, and committed artifacts are released
openly so that every number here is reproducible and every negative result is on the
record.

% =====================================================================
\section*{Data and code availability}
% =====================================================================

Code, configurations, provenance manifests, and the committed daily and
back-analysis artifacts (including \texttt{outlook-evidence-30d.json},
\texttt{benchmark.json}, \texttt{e12\_multiregion.json},
\texttt{e13\_catalog\_based.json}, and the per-region back-analyses) are at
\url{https://github.com/fsantibanezleal/CAOS_SEISMIC} under the MIT license; a
static forecast viewer is live at \url{https://seismic.fasl-work.com}. Catalog data
are from USGS ComCat (public domain) and the Centro Sismol\'ogico Nacional, Chile;
long-term anchoring from ISC-GEM; mechanisms from the Global CMT project; geodetic
velocities from the Nevada Geodetic Laboratory. Their licenses and required
attributions are tracked in the repository.

% =====================================================================
\begin{thebibliography}{99}
\itemsep 0pt

\bibitem{geller1997} R.~J. Geller, D.~D. Jackson, Y.~Y. Kagan, and F.~Mulargia,
``Earthquakes cannot be predicted,'' \emph{Science}, vol.~275, no.~5306, p.~1616,
1997, doi:10.1126/science.275.5306.1616.

\bibitem{jordan2011} T.~H. Jordan \emph{et al.}, ``Operational earthquake
forecasting: State of knowledge and guidelines for utilization,'' \emph{Ann.
Geophys.}, vol.~54, no.~4, 2011, doi:10.4401/ag-5350.

\bibitem{ogata1988} Y.~Ogata, ``Statistical models for earthquake occurrences and
residual analysis for point processes,'' \emph{J. Amer. Statist. Assoc.}, vol.~83,
no.~401, pp.~9--27, 1988, doi:10.1080/01621459.1988.10478560.

\bibitem{ogata1998} Y.~Ogata, ``Space--time point-process models for earthquake
occurrences,'' \emph{Ann. Inst. Statist. Math.}, vol.~50, no.~2, pp.~379--402,
1998, doi:10.1023/A:1003403601725.

\bibitem{reasenberg1989} P.~A. Reasenberg and L.~M. Jones, ``Earthquake hazard
after a mainshock in California,'' \emph{Science}, vol.~243, no.~4895,
pp.~1173--1176, 1989, doi:10.1126/science.243.4895.1173.

\bibitem{page2016} M.~T. Page, N.~van~der~Elst, J.~Hardebeck, K.~Felzer, and A.~J.
Michael, ``Three ingredients for improved global aftershock forecasts,'' \emph{Bull.
Seismol. Soc. Amer.}, vol.~106, no.~5, pp.~2290--2301, 2016, doi:10.1785/0120160073.

\bibitem{spassiani2023} I.~Spassiani, G.~Falcone, M.~Murru, and W.~Marzocchi,
``Operational earthquake forecasting in Italy: Validation after 10~yr of running,''
\emph{Geophys. J. Int.}, vol.~234, no.~3, pp.~2501--2518, 2023,
doi:10.1093/gji/ggad256.

\bibitem{serafini2025} F.~Serafini \emph{et al.}, ``A benchmark database of ten
years of prospective next-day earthquake forecasts,'' \emph{Sci. Data}, vol.~12,
p.~1501, 2025, doi:10.1038/s41597-025-05766-3.

\bibitem{stockman2026} S.~Stockman, D.~J. Lawson, and M.~J. Werner,
``EarthquakeNPP: Benchmark datasets for earthquake forecasting with neural point
processes,'' \emph{Trans. Mach. Learn. Res.}, 2026, arXiv:2410.08226.

\bibitem{devries2018} P.~M.~R. DeVries, F.~Vi\'egas, M.~Wattenberg, and B.~J.
Meade, ``Deep learning of aftershock patterns following large earthquakes,''
\emph{Nature}, vol.~560, pp.~632--634, 2018, doi:10.1038/s41586-018-0438-y.

\bibitem{mignan2019} A.~Mignan and M.~Broccardo, ``One neuron versus deep learning
in aftershock prediction,'' \emph{Nature}, vol.~574, pp.~E1--E3, 2019,
doi:10.1038/s41586-019-1582-8.

\bibitem{dascher2023} K.~Dascher-Cousineau, O.~Shchur, E.~E. Brodsky, and
S.~G\"unnemann, ``Using deep learning for flexible and scalable earthquake
forecasting,'' \emph{Geophys. Res. Lett.}, vol.~50, e2023GL103909, 2023,
doi:10.1029/2023GL103909.

\bibitem{zlydenko2023} O.~Zlydenko \emph{et al.}, ``A neural encoder for earthquake
rate forecasting,'' \emph{Sci. Rep.}, vol.~13, 2023, doi:10.1038/s41598-023-38033-9.

\bibitem{bird2015} P.~Bird, D.~D. Jackson, Y.~Y. Kagan, C.~Kreemer, and R.~S.
Stein, ``GEAR1: A global earthquake activity rate model constructed from geodetic
strain rates and smoothed seismicity,'' \emph{Bull. Seismol. Soc. Amer.}, vol.~105,
no.~5, pp.~2538--2554, 2015, doi:10.1785/0120150058.

\bibitem{bayona2022} J.~A. Bayona \emph{et al.}, ``Prospective evaluation of
multiplicative hybrid earthquake forecasting models in California,'' \emph{Geophys.
J. Int.}, vol.~229, no.~3, pp.~1736--1753, 2022, doi:10.1093/gji/ggac018.

\bibitem{scordilis2006} E.~M. Scordilis, ``Empirical global relations converting
$M_S$ and $m_b$ to moment magnitude,'' \emph{J. Seismol.}, vol.~10, pp.~225--236,
2006, doi:10.1007/s10950-006-9012-4.

\bibitem{wiemer2000} S.~Wiemer and M.~Wyss, ``Minimum magnitude of completeness in
earthquake catalogs,'' \emph{Bull. Seismol. Soc. Amer.}, vol.~90, no.~4,
pp.~859--869, 2000, doi:10.1785/0119990114.

\bibitem{woessner2005} J.~Woessner and S.~Wiemer, ``Assessing the quality of
earthquake catalogues: Estimating the magnitude of completeness and its
uncertainty,'' \emph{Bull. Seismol. Soc. Amer.}, vol.~95, no.~2, pp.~684--698,
2005, doi:10.1785/0120040007.

\bibitem{zhuang2002} J.~Zhuang, Y.~Ogata, and D.~Vere-Jones, ``Stochastic
declustering of space--time earthquake occurrences,'' \emph{J. Amer. Statist.
Assoc.}, vol.~97, no.~458, pp.~369--380, 2002, doi:10.1198/016214502760046925.

\bibitem{zaliapin2013} I.~Zaliapin and Y.~Ben-Zion, ``Earthquake clusters in
southern California I: Identification and stability,'' \emph{J. Geophys. Res. Solid
Earth}, vol.~118, no.~6, pp.~2847--2864, 2013, doi:10.1002/jgrb.50179.

\bibitem{blewitt2016} G.~Blewitt, C.~Kreemer, W.~C. Hammond, and J.~Gazeaux,
``MIDAS robust trend estimator for accurate GPS station velocities without step
detection,'' \emph{J. Geophys. Res. Solid Earth}, vol.~121, no.~3, pp.~2054--2068,
2016, doi:10.1002/2015JB012552.

\bibitem{bird2003} P.~Bird, ``An updated digital model of plate boundaries,''
\emph{Geochem. Geophys. Geosyst.}, vol.~4, no.~3, 1027, 2003,
doi:10.1029/2001GC000252.

\bibitem{helmstetter2007} A.~Helmstetter, Y.~Y. Kagan, and D.~D. Jackson,
``High-resolution time-independent grid-based forecast for $M \ge 5$ earthquakes in
California,'' \emph{Seismol. Res. Lett.}, vol.~78, no.~1, pp.~78--86, 2007,
doi:10.1785/gssrl.78.1.78.

\bibitem{savran2022} W.~H. Savran \emph{et al.}, ``pyCSEP: A Python toolkit for
earthquake forecast developers,'' \emph{Seismol. Res. Lett.}, vol.~93, no.~5,
pp.~2858--2870, 2022, doi:10.1785/0220220033.

\bibitem{savran2020} W.~H. Savran, M.~J. Werner, W.~Marzocchi, \emph{et al.},
``Pseudoprospective evaluation of UCERF3-ETAS forecasts during the 2019 Ridgecrest
sequence,'' \emph{Bull. Seismol. Soc. Amer.}, vol.~110, no.~4, pp.~1799--1817,
2020, doi:10.1785/0120200026.

\bibitem{schorlemmer2007} D.~Schorlemmer, M.~C. Gerstenberger, S.~Wiemer, D.~D.
Jackson, and D.~A. Rhoades, ``Earthquake likelihood model testing,'' \emph{Seismol.
Res. Lett.}, vol.~78, no.~1, pp.~17--29, 2007, doi:10.1785/gssrl.78.1.17.

\bibitem{zechar2010} J.~D. Zechar, M.~C. Gerstenberger, and D.~A. Rhoades,
``Likelihood-based tests for evaluating space--rate--magnitude earthquake
forecasts,'' \emph{Bull. Seismol. Soc. Amer.}, vol.~100, no.~3, pp.~1184--1195,
2010, doi:10.1785/0120090192.

\bibitem{rhoades2011} D.~A. Rhoades \emph{et al.}, ``Efficient testing of
earthquake forecasting models,'' \emph{Acta Geophys.}, vol.~59, no.~4,
pp.~728--747, 2011, doi:10.2478/s11600-011-0013-5.

\bibitem{kagan2017} Y.~Y. Kagan, ``Earthquake number forecasts testing,''
\emph{Geophys. J. Int.}, vol.~211, no.~1, pp.~335--345, 2017, doi:10.1093/gji/ggx300.

\bibitem{marzocchi2012} W.~Marzocchi, J.~D. Zechar, and T.~H. Jordan, ``Bayesian
forecast evaluation and ensemble earthquake forecasting,'' \emph{Bull. Seismol.
Soc. Amer.}, vol.~102, no.~6, pp.~2574--2584, 2012, doi:10.1785/0120110327.

\bibitem{herrmann2023} M.~Herrmann and W.~Marzocchi, ``Maximizing the forecasting
skill of an ensemble model,'' \emph{Geophys. J. Int.}, vol.~234, no.~1,
pp.~73--87, 2023, doi:10.1093/gji/ggad020.

\bibitem{gneiting2007} T.~Gneiting and A.~E. Raftery, ``Strictly proper scoring
rules, prediction, and estimation,'' \emph{J. Amer. Statist. Assoc.}, vol.~102,
no.~477, pp.~359--378, 2007, doi:10.1198/016214506000001437.

\bibitem{zechar2008} J.~D. Zechar and T.~H. Jordan, ``Testing alarm-based
earthquake predictions,'' \emph{Geophys. J. Int.}, vol.~172, no.~2, pp.~715--724,
2008, doi:10.1111/j.1365-246X.2007.03676.x.

\end{thebibliography}

\end{document}
